\documentclass[12pt,a4paper]{article}

% Packages
\usepackage[utf8]{inputenc}
\usepackage[T1]{fontenc}
\usepackage{graphicx}
\usepackage{amsmath,amssymb}
\usepackage{booktabs}
\usepackage{caption}
\usepackage{subcaption}
\usepackage{float}
\usepackage{hyperref}
\usepackage{geometry}

\geometry{
    left=1in,
    right=1in,
    top=1in,
    bottom=1in
}

\title{ECE131a Extra Credit Project\\ {Distributed Detection in Gaussian Noise}}
\author{Wesley Hon}
\date{\today}

\begin{document}

%-------------------------------------------------
% Title Page
%-------------------------------------------------
\maketitle

%-------------------------------------------------
% Front Matter
%-------------------------------------------------
\pagenumbering{roman}

\tableofcontents
\clearpage

%-------------------------------------------------
% Main Matter
%-------------------------------------------------
\pagenumbering{arabic}

\section{Introduction}

\subsection{Background}
This problem is adapted from the results in [1] and [2]. A scalar signal
$S \in \{-m,m\}$ is transmitted from a satellite to a field of $n$ sensors
on the surface of the earth. $S$ communicates an important message, and is
equally likely to take either value so that

\[
P_S(s)=
\begin{cases}
\frac{1}{2}, & \text{for } s=-m,\\[4pt]
\frac{1}{2}, & \text{for } s=m,\\[4pt]
0, & \text{otherwise}.
\end{cases}
\]

Sensor $i$ observes random variable

\[
X_i = S + Z_i
\]

where the random variables $Z_i$ are independent and identically distributed
according to a standard normal distribution $\mathcal{N}(0,1)$.
This project explores how to efficiently and reliably estimate $S$ from the
vector of observations

\[
\mathbf{X} = [X_1, X_2, \ldots, X_n].
\]

Any estimate $\hat{S}$ will naturally depend on the observables.
A \emph{statistic} is a function of the set of observables that we use in
preparing our estimate. For this problem we will be especially interested in
the statistic

\[
T_n = \frac{1}{n}\sum_{i=1}^{n} X_i,
\]

which is sometimes called the sample mean.

A \emph{sufficient statistic} of a set of observations for estimation of a
particular random variable is one that allows as good of an estimate of the
random variable as the original observations themselves.
$T_n$ is a sufficient statistic of $\mathbf{X}$ for $S$ if and only if

\[
f_{X|T,S}(x|t,s) = f_{X|T}(x|t).
\]

In our specific case, this turns out to be true so that the sample mean is a
sufficient statistic of $\mathbf{X}$ for $S$. We can use the sample mean
$T_n$ to estimate $S$ and be confident that we have not lost any information.
Specifically, we will estimate $S$ as follows:

\[
\hat{S} =
\begin{cases}
m, & \text{if } T_n \ge 0,\\[4pt]
-m, & \text{otherwise}.
\end{cases}
\]

\subsection{Report Organization}
The project is 25 project points which is equivalent to 5\% of the total grade percentage as extra credit.   The project must be completed to earn an A+ grade.\\
There are 8 problems, each answer must have a Matlab plot to back up the answer. \\

Project Statement:\\
\noindent \href{https://drive.google.com/drive/folders/1uo9Ji4P91C4RJN64mATiQfMw05c1TiUV}{ECE 131a Spring 2026 Extra Credit Project, Distributed Detection in Gaussian Noise}

%-------------------------------------------------

\section{Problem 1}
3pts.\\

\hrule 

\

\noindent Suppose that $m = 0.5$. Use Matlab function \texttt{semilogy} to plot the probability of error $P(\hat{S} \neq S)$ as a function of the number of sensors $n$ for $1 \leq n \leq 100$. The Matlab function \texttt{qfunc} is your friend for this problem. How many sensors do you need so that $P(\hat{S} \neq S) < 10^{-3}$? Call this number $n_g$ (as in the number of Gaussian samples).\\

\hrule

%-------------------------------------------------


\end{document}